
% SGA 3，报告 VII_B；英文重建本第一轮的中文译本。
% 范围：依据文献局部页 57--58，第 2.5.2 节及注 2.5.2.A 全文，止于 2.5.3 前。
% 合并阅读本页 570--571；重版印刷页 559--560；原版边码 524；编者注 100--102。
% 依据：Polo--Gille Exp7B-13oct24.pdf。
% SHA-256 D303690E0195654BBE7C3DAA5D89EA23E24CD709A3575217234AC19AE8836694。
% 参照了 Jacob Reinhold 的 CC BY 4.0 Markdown，作为对照来源。
% 两幅图均为第一轮直接从 PDF 截取的原图。
\subsubsection*{2.5.2}
\phantomsection
\label{sga3:VIIB:2.5.2}
设 \(k\) 是完美域。在这种情形，\SGARefLink{sga3:VIIB:1.6.2}{1.6.2} 已说明态射
\[
  p:X\longmapsto X_e
\]
有截面
\[
  s:X_e\longrightarrow X
\]
且该截面关于 \(X\) 是函子性的。因此，当 \(X\) 是形式群时，这个截面是形式群态射。于是得到：
\medskip
\noindent\textbf{命题。}
\phantomsection
\label{sga3:VIIB:2.5.2:proposition}
{\itshape 当 \(k\) 是完美域时，每个形式 \(k\)-群 \(G\) 都典范同构于无穷小群 \(G_0\) 与作用于 \(G_0\) 的平展群 \(G_e\) 的半直积。}\footnote{%
\label{sga3:VIIB:nde:100}同一论证还表明，对任意域 \(k\)，若 \(G\) 的所有剩余域都等于 \(k\)，则 \(G_e\) 是常值 \(k\)-群 \((M)_k\)，其中
\[
  M=G(k)=\operatorname{Spf}^{*}\mathbf H(G),
\]
且 \(\mathbf H(G)\) 是 \(\mathbf H(G_0)\) 与群代数 \(kM\) 的半直积；参见下文 \SGARefLink{sga3:VIIB:2.9}{2.9} 的补充。}
若 \(G\) 还是交换的，则 \(G\) 是 \(G_0\) 与 \(G_e\) 的积。由 \SGARefLink{sga3:VIIB:2.2.2}{2.2.2}，交换形式 \(k\)-群的这个典范分解，对应于仿射交换 \(k\)-群概形的类似分解；参见下文 \SGARefLink{sga3:VIIB:2.5.3}{2.5.3}。
\medskip
\noindent\textbf{注 2.5.2.A。}
\phantomsection
\label{sga3:VIIB:2.5.2.A}
\footnote{\label{sga3:VIIB:nde:101}此注由编者添加。}
当 \(k\) 不是完美域时，正合列
\[
  1\longrightarrow G_0\longrightarrow G\longrightarrow G_e
  \longrightarrow 1
\]
不一定分裂。考察下面取自 \([DG70]\)，§III.6，8.6 的例子。设 \(k\) 是特征 \(p>0\) 的非完美域，选取
\[
  \lambda\in k\setminus k^p,
\]
并令 \(L_\lambda\) 为具有基 \((x,y)\) 的阿贝尔限制 \(p\)-Lie 代数，其符号 \(p\) 次幂\textup{（参见 VII\textsubscript{A}，\SGARefLink{sga3:VIIA:5.2}{5.2}）}由下式给定
\[
  x^{(p)}=x,
  \qquad
  y^{(p)}=\lambda x.
\]
令
\[
  U_p(L_\lambda)
  =
  k[x,y]/(x^p-x,\ y^p-\lambda x)
\]
为 \(L_\lambda\) 的限制包络代数\textup{（参见 VII\textsubscript{A}，\SGARefLink{sga3:VIIA:5.3}{5.3}）}，并令 \(G_\lambda\) 为以 \(U_p(L_\lambda)\) 为仿射代数的交换形式 \(k\)-群。
则 \(G_\lambda\) 是常值平展 \(k\)-群 \((\mathbb Z/p\mathbb Z)_k\) 按无穷小 \(k\)-群
\[
  \alpha_{p,k}=\operatorname{Spf}k[x]/(x^p).
\]
的一个非分裂扩张。
\phantomsection
\label{sga3:VIIB:2.5.2.A:alpha-p-source-note}
\noindent\textit{原文说明。}\enspace
法文重版此处印作 \(\operatorname{Spf}k[x]/(x^p-x)\)。这个代数是平展的，而所断言的核是标准无穷小群 \(\alpha_{p,k}=\operatorname{Spf}k[x]/(x^p)\)。事实上，\(U_p(L_\lambda)\) 中核纤维 \(x=0\) 的坐标代数是 \(k[y]/(y^p)\)。英文本采用标准无穷小群的读法，并在此注中保留印本读法。
于是有交换形式 \(k\)-群的非分裂正合列：\footnote{%
\label{sga3:VIIB:nde:102}群 \(\alpha_{p,k}\)、\((\mathbb Z/p\mathbb Z)_k\) 和 \(G_\lambda\) 也都是有限平坦代数 \(k\)-群；参见 \SGARefLink{sga3:VIIB:2.2.2}{2.2.2} 的编者注 84。}
\begin{center}
  \phantomsection
  \label{sga3:VIIB:2.5.2.A:nonsplit-formal-sequence}
  % SGA3-VII-DIAG-116
  % SGA3-VII-DIAG-116
% Authority: Exp7B-13oct24.pdf, local p. 58.
% Crop: figures/Exp7B_localp058_nonsplit_formal_group_exact_sequence.png
% Locator: tex/components/49_expose_VIIB_section252_perfect_field_splitting_en.tex, Loop 1 line 113.
\(
0\longrightarrow\boldsymbol{\alpha}_{p,k}
\longrightarrow G_\lambda
\longrightarrow(\mathbb Z/p\mathbb Z)_k
\longrightarrow0.
\)

\end{center}
通过对偶，它对应于交换代数 \(k\)-群的非分裂正合列：
\begin{center}
  \phantomsection
  \label{sga3:VIIB:2.5.2.A:nonsplit-dual-sequence}
  % SGA3-VII-DIAG-117
  % SGA3-VII-DIAG-117
% Authority: Exp7B-13oct24.pdf, local p. 58.
% Crop: figures/Exp7B_localp058_nonsplit_dual_algebraic_group_exact_sequence.png
% Locator: tex/components/49_expose_VIIB_section252_perfect_field_splitting_en.tex, Loop 1 line 122.
\(
0\longrightarrow\boldsymbol{\mu}_{p,k}
\longrightarrow D'(G_\lambda)
\longrightarrow\boldsymbol{\alpha}_{p,k}
\longrightarrow0
\)

\end{center}
其中
\[
  \mu_{p,k}=\operatorname{Spec}k[t]/(t^p-1).
\]
用这种方法可得到 \(\alpha_{p,k}\) 按 \(\mu_{p,k}\) 的所有扩张；参见 \([DG70]\)，§III.6，8.6。
